\clearpage
\part{第三题：字符串解析与综合模拟}
\noindent 本篇围绕长题意的实现展开：先解析输入，再组织状态，最后执行规则并检查边界。遇到陌生背景时，先列出输入格式、状态定义、修改规则和查询规则，再写程序。

\section{字符串处理与解析基础}
\subsection{常用操作与区间含义}
\ojpractice{CSP201409C}{字符串匹配：第2次认证，按是否区分大小写查找子串，训练扫描与比较}
\ojpractice{CSP202409B}{字符串变换：第35次认证，按规则逐轮替换，完整解答见第二篇练习题}
string 的下标从0开始。\texttt{s.substr(pos,len)} 的第二个参数是长度；\texttt{s.find(t)} 返回首次匹配位置，找不到时返回 \texttt{string::npos}。不要把 npos 随意存进 int 后再与某个猜测值比较。字符串比较是字典序，例如 \texttt{"10" < "2"}，与整数大小不同。

扫描字符串时最好约定：“指针 p 总是指向下一个尚未处理的字符”。每个分支要么推进 p，要么返回。若某个分支既不推进也不退出，就容易死循环。处理完整段落时，通常比反复在字符串中间 erase 更清楚、也更高效。

\example{带引号的单行分词}
\textbf{题目。}一行文本由普通字符、空格和英文双引号组成。引号外的连续空格分隔单词；引号内的空格属于单词；双引号本身不输出。相邻的普通片段和引号片段属于同一个单词。空引号表示一个空单词。本题不支持转义字符。若引号不匹配，输出 ERROR。行长不超过 $10^5$。

\textbf{输入。}一整行文本，可为空。\textbf{输出。}合法时先输出单词数，随后每个单词单独一行，用方括号包住，便于看见空单词和空格。
\begin{verbatim}
输入
add "hello world" "" x
输出
4
[add]
[hello world]
[]
[x]
\end{verbatim}
\textbf{状态设计。}quoted 表示当前是否在引号内；word[0..len-1] 保存当前单词；active 表示当前单词已经开始。active 不能由 len 是否为0代替，否则空引号会丢失。遇到引号切换 quoted；遇到引号外空格时结束当前单词；其余字符追加到word并增加len。扫描结束后检查引号是否闭合，再处理最后一个单词。

\textbf{正确性。}每个字符恰好进入一种分支。引号仅改变状态，分隔空格仅结束一个已开始的单词，因此保留的字符、单词边界和空单词均符合规则。时间 $O(L)$，输出存储 $O(L)$。
\Code{04_tokens.cpp}

\subsection{C字符数组与string的区别}
C 字符串以终止字符 \verb|'\0'| 结尾，数组需要额外空间存放它。strlen 统计终止符前的字符数。定义后的字符数组不能用赋值表达式整体赋值；string 可以直接赋值。strcmp 返回值的符号表示字典序关系，不保证只返回1或-1。字符分类函数如 isdigit 接收字符时，可先转为 unsigned char，避免负 char 值带来的未定义行为。

\subsection{解析与业务处理分开}
例如命令 \texttt{GRANT alice file1}，解析层先产生三个字符串，再由执行层决定如何修改权限。这样出现错误时，可以先打印解析结果，判断错在读入还是状态更新。本篇下一节会进一步用函数处理路径和表达式。

\exercise
\textbf{练习1。}输入 \verb|a"b c"d|，分词结果是什么？\textbf{解答。}只有一个单词 \texttt{ab cd}。相邻片段没有引号外空格，不应拆开。

\textbf{练习2。}输入只有三个空格，输出什么？\textbf{解答。}单词数为0。输入 \verb|""| 时则是一个空单词，输出1和 \texttt{[]}。

\textbf{练习3。}若题目新增反斜杠转义，当前程序能直接使用吗？\textbf{解答。}不能。需要新增 escape 状态：读取反斜杠后，下一个字符按普通字符处理，不切换引号状态。还要定义行末孤立反斜杠是否非法。规则改变后必须重新列出状态转移。
\newpage
